\documentclass[../DoAn.tex]{subfiles}
\begin{document}

\begin{center}
    \Large{\textbf{TÓM TẮT NỘI DUNG ĐỒ ÁN}}\\
\end{center}
\vspace{1cm}
Sinh viên viết tóm tắt ĐATN của mình trong mục này, với 200 đến 350 từ. Theo trình tự, các nội dung tóm tắt cần có: (i) Giới thiệu vấn đề (tại sao có vấn đề đó, hiện tại được giải quyết chưa, có những hướng tiếp cận nào, các hướng này giải quyết như thế nào, hạn chế là gì), (ii) Hướng tiếp cận sinh viên lựa chọn là gì, vì sao chọn hướng đó, (iii) Tổng quan giải pháp của sinh viên theo hướng tiếp cận đã chọn, và (iv) Đóng góp chính của ĐATN là gì, kết quả đạt được sau cùng là gì. Sinh viên cần viết thành đoạn văn, không được viết ý hoặc gạch đầu dòng.

Phần tóm tắt cần nêu được các ý: vấn đề cần thực hiện; phương pháp thực hiện; công cụ sử dụng (phần mềm, phần cứng…); kết quả của đồ án có phù hợp với các vấn đề đã đặt ra hay không; tính thực tế của đồ án, định hướng phát triển mở rộng của đồ án (nếu có); các kiến thức và kỹ năng mà sinh viên đã đạt được.

(Đối với luận văn thạc sĩ, phần tóm tắt được in trong một bản nộp riêng)
\begin{flushright}
Sinh viên thực hiện\\
\begin{tabular}{@{}c@{}}
\textit{(Ký và ghi rõ họ tên)}
\end{tabular}
\end{flushright}

\end{document}